\section{MIKU methodology}
\label{sec:method}

\subsection{Overview}

MIKU reconstructs the two pieces of information lost by time-blind and obstacle-wise PVD. It first estimates the ego arrival time $\tau(s)$ and queries time-dependent obstacle occupancy. It then computes threat-dependent center bounds, groups longitudinally connected obstacles, and selects continuous lateral bands. After the path QP, uncertainty-expanded occupancy tubes define safe passage windows. A layered temporal homotopy graph selects a causally compatible window sequence and writes it into the ST position bounds. Figure~\ref{fig:miku_method_flow_en} summarizes the data flow.

\begin{figure}
\centering
\adjustbox{max width=\textwidth}{\begin{tikzpicture}[x=.92cm,y=.92cm,>=Stealth]
\node[modin,minimum width=2.55cm] (input) at (-0.5,1.05) {Road geometry\\Ego state\\Obstacle predictions};
\node[modimp,minimum width=2.55cm] (occupancy) at (3.35,1.05) {Arrival-time query\\Threat margin\\Robust occupancy};
\node[modimp,minimum width=2.45cm,text width=2.35cm] (spatial) at (6.70,1.05) {Continuous partition\\Top-$K$ spatial classes};
\node[mod,minimum width=1.95cm] (path) at (9.55,1.05) {Path QP\\$l(s)$};
\node[modimp,minimum width=2.45cm,text width=2.35cm] (temporal) at (6.70,-1.05) {Safe windows\\Temporal homotopy};
\node[mod,minimum width=1.95cm] (speed) at (9.55,-1.05) {Speed QP\\$s(t)$};
\node[modout,minimum width=2.35cm] (out) at (12.35,-1.05) {Space--time\\trajectory};
\draw[flow] (input)--(occupancy);
\draw[flowimp] (occupancy)--(spatial);
\draw[flowimp] (spatial)--(path);
\draw[flowimp] (spatial)--(temporal);
\draw[flow] (path.south) to[out=-160,in=90] (temporal.north east);
\draw[flowimp] (temporal)--(speed);
\draw[flow] (speed)--(out);
\draw[flow,dashed] (speed.south) .. controls (9.35,-2.35) and (3.35,-2.35) ..
  node[below,font=\scriptsize]{arrival-time update} (occupancy.south);
% \node[font=\scriptsize,align=center] at (5.00,2.08) {MIKU constraint compiler};
% \node[font=\scriptsize,align=center] at (10.70,-2.78) {Existing convex solvers retained};
\end{tikzpicture}
}
\caption{MIKU constraint construction and transfer through a PVD planning stack.}
\label{fig:miku_method_flow_en}
\end{figure}

\subsection{Threat-aware margins}
\label{sec:method_threat_en}

For obstacle $i$, MIKU combines collision urgency, lateral overlap, relative motion, object type, and local interaction density into a normalized threat score:
\begin{equation}
 \Theta_i=w_1f_{\mathrm{TTC}}(i)+w_2f_{\mathrm{overlap}}(i)+w_3f_{\mathrm{vel}}(i)+w_4f_{\mathrm{type}}(i)+w_5f_{\mathrm{inter}}(i).
 \label{eq:threat_score_en}
\end{equation}
All factors lie in $[0,1]$, and the weights are $(0.30,0.20,0.15,0.10,0.25)$. When the ego vehicle is behind an obstacle and closing, $\mathrm{TTC}_i=(s_i-s_{\mathrm{ego}})/(\dot{s}_{\mathrm{ego}}-\dot{s}_i)$; otherwise it is $+\infty$. The piecewise-linear TTC factor equals one for $\mathrm{TTC}_i\leq2.0$ s, zero for $\mathrm{TTC}_i\geq7.0$ s, and decreases linearly in between \citep{minderhoud2001extended}. The overlap and relative-speed factors are
\begin{align}
 f_{\mathrm{overlap}}(i)&=\min\!\left(1,\frac{\max(0,\min(l_i^+,l_{\mathrm{ego}}^+)-\max(l_i^-,l_{\mathrm{ego}}^-))}{\max(0.1,\min(W_i,W_{\mathrm{ego}}))}\right),\\
 f_{\mathrm{vel}}(i)&=\sigma(v_{\mathrm{rel},i}/v_{\max}),\qquad \sigma(x)=\frac{1}{1+e^{-5x}}.
\end{align}
Here $W_i=l_i^+-l_i^-$ is the obstacle width; the cap keeps the overlap factor in $[0,1]$ for obstacles wider than the ego vehicle.
The type factors for vulnerable road users, motor vehicles, unknown movable objects, static obstacles, and small markers are respectively $1.0$, $0.7$, $0.5$, $0.3$, and $0.15$. Density uses a 10 m neighborhood and is zero for a single obstacle. The extra margin is
\begin{equation}
 d_{\mathrm{buf},i}=d_{\mathrm{buf,min}}+(d_{\mathrm{buf,max}}-d_{\mathrm{buf,min}})\Theta_i,
 \label{eq:dynamic_delta_en}
\end{equation}
where $[d_{\mathrm{buf,min}},d_{\mathrm{buf,max}}]=[0.10,0.40]$ m. This term changes the boundary value only; it does not change the sorting or scan structure.

\subsection{Continuous lateral partition and maximum gap}
\label{sec:method_band_en}

Obstacles whose longitudinal intervals overlap are assigned to the same connected component. Sorting by the lower station edge and maintaining the largest upper edge constructs all components in $O(n\log n)$. Within one component, sort the active intervals by $u_i$ (ties by descending $v_i$). The following lemma limits the side-assignment search.

\begin{lemma}[Continuous optimal partition]
\label{lem:continuous_partition_en}
For the center-band objective in (\ref{eq:center_band_en}), there is an optimal assignment in which the obstacles sorted by $u_i$ form a prefix assigned to $L$ and a suffix assigned to $R$. Hence only $k+1$ split points need be considered.
\end{lemma}

\emph{Proof.} If both sides are nonempty, let $\beta=\min_{i\in\mathcal{R}}u_i$. Move every obstacle currently in $\mathcal{L}$ with $u_i\geq\beta$ to $\mathcal{R}$. The upper bound remains unchanged because each moved lower edge is no smaller than $\beta$, whereas the lower bound cannot increase after removing elements from $\mathcal{L}$. The width therefore does not decrease. Repeating this operation yields a continuous split. Empty-side cases give $p=0$ or $p=k$. Descending $v_i$ resolves ties in $u_i$. \hfill$\Box$

For split $p$, define the prefix maximum and candidate upper edge by
\begin{equation}
 V_0=L^-,\qquad V_p=\max\left(L^-,\max_{1\leq q\leq p}v_{(q)}\right),
 \label{eq:prefix_v_en}
\end{equation}
\begin{equation}
 U_p=\begin{cases}\min(L^+,u_{(p+1)}),&p<k,\\L^+,&p=k,\end{cases}\qquad g_p=U_p-V_p.
 \label{eq:gap_def_en}
\end{equation}

\begin{theorem}[Maximum-gap optimality]
\label{thm:max_gap_en}
For the fixed-section interval model, the optimal center-band width is $W^*=\max_{p=0,\ldots,k}g_p$. The candidates are obtained by sorting and one prefix scan, with complexity $O(k\log k)$.
\end{theorem}

\emph{Proof.} For split $p$, (\ref{eq:center_band_en}) gives the lower edge $V_p$ and the tightest upper edge $U_p$, so its width is $g_p$. Lemma~\ref{lem:continuous_partition_en} guarantees an optimal continuous split; enumerating the $k+1$ splits therefore attains the optimum. \hfill$\Box$

The theorem is exact only for this fixed-section interval abstraction. It does not certify a continuous global optimum for the full nonconvex trajectory problem. MIKU keeps the $Q=3$ widest positive bands in each component and uses dynamic programming to balance band width against lateral-center changes. If band $B_{jq}$ has center $c_{jq}$ and width $g_{jq}$ in component $j$, the sequence minimizes
\begin{equation}
 -\sum_jg_{jq_j}+\lambda_s\left(|c_{1q_1}-l_0|+\sum_{j>1}|c_{jq_j}-c_{j-1,q_{j-1}}|\right),
 \label{eq:spatial_homotopy_en}
\end{equation}
with the recurrence
\begin{equation}
 D_{jq}=-g_{jq}+\min_r\{D_{j-1,r}+\lambda_s|c_{jq}-c_{j-1,r}|\}.
 \label{eq:spatial_dp_en}
\end{equation}

\begin{figure}
\centering
\adjustbox{max width=0.82\textwidth}{\begin{tikzpicture}[x=1cm,y=1cm]
  \fill[bglight] (0,-.45) rectangle (10,.45);
  \draw[gray,thick] (0,0)--(10,0);
  \foreach \x/\name in {1.2/$O_1$,3.0/$O_2$,6.0/$O_3$,8.1/$O_4$}{
    \fill[dangerred!70] (\x-.38,-.25) rectangle (\x+.38,.25);
    \node[font=\scriptsize] at (\x,.7) {\name};
  }
  \draw[deepblue,very thick,<->] (3.38,-.85)--(5.62,-.85);
  \node[deepblue,font=\scriptsize] at (4.5,-1.15) {$g_2=U_2-V_2$ (maximum gap)};
  \foreach \x/\lab in {0/$p=0$,2.62/$p=1$,4.62/$p=2$,6.62/$p=3$,8.48/$p=4$}{
    \draw[gray!70,dashed] (\x,-.48)--(\x,.48);
    \node[font=\tiny] at (\x,-.63) {\lab};
  }
  \node[font=\scriptsize,align=center] at (5,1.15) {$k=4$ intervals yield $k+1=5$ continuous split candidates};
\end{tikzpicture}
}
\caption{Continuous split candidates for the fixed-section maximum-gap problem. Four sorted forbidden intervals induce five candidate splits; the prefix maximum is required when interval upper edges are nested.}
\label{fig:max_gap_en}
\end{figure}

\subsection{Arrival-time-dependent projection}
\label{sec:method_corridor_en}

At station $s$, the path stage queries obstacle occupancy at the estimated ego arrival time. With the current state and a constant-acceleration approximation,
\begin{equation}
 \tau(s)=\frac{-v_{\mathrm{ego}}+\sqrt{v_{\mathrm{ego}}^2+2a_{\mathrm{ego}}(s-s_{\mathrm{ego}})}}{a_{\mathrm{ego}}},
 \label{eq:arrival_time_en}
\end{equation}
with the constant-speed limit used when $|a_{\mathrm{ego}}|<10^{-3}$. A negative discriminant is marked outside the planning horizon. In rolling replanning, $\tau(s)$ is interpolated from the previous speed solution and updated at each cycle.

The numerical scenarios use $s_i(t)=s_{i,0}+v_{s,i}t$ and $l_i(t)=l_{i,0}+v_{l,i}t$. For bounded prediction errors, the longitudinal and lateral expansion at time $t$ are
\begin{equation}
 \epsilon_{s,i}(t)=\epsilon_{s0,i}+t\epsilon_{vs,i},\qquad
 \epsilon_{l,i}(t)=\epsilon_{l0,i}+t\epsilon_{vl,i}.
 \label{eq:prediction_tube_en}
\end{equation}
The physical obstacle and ego footprint are expanded by these quantities before the ST occupancy map is formed. Nominal $\tau(s)$ selects and ranks alternatives; collision checking uses the expanded occupancy.

\subsection{Temporal homotopy and ST transfer}
\label{sec:method_inject_en}

For conflict region $c$, let $I_{cr}=[t_{cr}^{in},t_{cr}^{out}]$ be each overlap interval between the path and the predicted occupancy tube. After temporal safety expansion $\eta_c$, the safe windows are
\begin{equation}
 \mathcal{W}_c=[0,T]\setminus\bigcup_r[t_{cr}^{in}-\eta_c,t_{cr}^{out}+\eta_c].
 \label{eq:safe_windows_en}
\end{equation}
Intersecting these windows with the earliest reachable arrival time yields $W_{cq}=[\alpha_{cq},\beta_{cq}]$. For conflict points ordered by station, connect adjacent nodes when
\begin{equation}
 \hat t_c\geq\hat t_{c-1}+\frac{s_c-s_{c-1}}{v_{\max}},\qquad \hat t_c\in W_{cq}.
 \label{eq:temporal_edge_en}
\end{equation}
Node cost measures deviation from $\tau(s_c)$, and edge cost discourages unnecessary switching between before- and after-obstacle passage. A bounded-width dynamic program selects one causal sequence.

For a selected window, its semantics enter the speed QP as linear box constraints:
\begin{equation}
 \begin{aligned}
 s_j^{ub}&=\min(s_j^{ub,0},s_c^- -\epsilon_s),&&t_j<\alpha_{cq},\\
 s_j^{lb}&=\max(s_j^{lb,0},s_c^+ +\epsilon_s),&&t_j\geq\beta_{cq}.
 \end{aligned}
 \label{eq:window_to_st_en}
\end{equation}
If no complete window sequence is reachable, the planner imposes a stopping upper bound before the first unreachable conflict. The downstream speed QP remains a convex QP with the same objective and kinematic equalities.

\subsection{Algorithm and complexity}

\begin{algorithm}
\small
\caption{MIKU constraint compilation}
\label{alg:miku_en}
\KwIn{Road bounds, ego state, and predicted trajectories for $n$ obstacles}
\KwOut{Center path bounds and ST position bounds}
Estimate $\tau(s)$ and compute threat-dependent center intervals for each obstacle.\;
Sort station intervals and form longitudinal connected components.\;
\ForEach{component}{
  Sort by $(u_i,-v_i)$, compute prefix maxima, and retain the $Q$ widest positive bands.\;
}
Select a continuous spatial-homotopy sequence by (\ref{eq:spatial_homotopy_en}).\;
Write the selected side assignments into the active path bounds at each station.\;
Construct uncertainty-expanded occupancy tubes and safe windows.\;
Select a causal temporal-homotopy sequence by (\ref{eq:temporal_edge_en}) and compile it into (\ref{eq:window_to_st_en}).\;
If enabled, update $\tau(s)$ from the speed solution and repeat until the iteration limit or convergence.\;
\Return{Path and ST bounds}
\end{algorithm}

For $K$ path samples, $m$ components, at most $Q$ spatial candidates per component, and a bounded temporal beam width $B$, sorting and grouping cost $O(n\log n)$. With $K_s$ retained prefixes per terminal spatial band, the layered spatial dynamic program costs $O(mK_sQ^2\log(K_sQ))$; temporal beam expansion and ranking cost $O(hBR\log(BR))$ for $h$ conflict points and $R$ reachable windows per layer. The implementation uses $Q=3$, $K_s=3$, and $B=8$. Occupancy queries are $O(nK)$ and naive interaction-density computation is $O(n^2)$ in the worst case.
